\documentclass[11pt,a4paper]{article}
\usepackage[T1]{fontenc}
\usepackage{lmodern,amsmath,amssymb,booktabs,microtype}
\usepackage[margin=24mm]{geometry}
\usepackage[colorlinks=true,urlcolor=blue,linkcolor=blue]{hyperref}
\newcommand{\EE}{\mathbb E}
\newcommand{\trm}{\operatorname{tr}_m}
\DeclareMathOperator{\tr}{tr}
\newcommand{\Qc}{\mathcal Q}
\newcommand{\Pc}{\mathcal P}
\title{The temporal singlet as a compressed evolution\\An exact memory kernel and its fluctuation seed}
\author{Alan Day \and Chengheng\thanks{AI research assistant (OpenAI).}}
\date{Local research note --- 6 November 2026}
\begin{document}\maketitle
\begin{abstract}
The full temporal singlet admits a finite-system isometric dilation
through the adjacent charge sector. This gives its positive kernel,
an exact composition defect and a closed equation with memory, without
replacing unequal-time fermion anticommutators by scalars. The initial
memory operator contains the previously computed dynamical singlet.
A Duhamel bound proves that the ordinarily normalized bilocal tends
to the identity uniformly on bounded string-time intervals in double
scaling. The physical reconstruction problem therefore concerns its
renormalized fluctuations and the full orthogonal memory dynamics.
\end{abstract}
\section{A charge-removal isometry}
Use the standard complex CAR and the neutral $p$-body Hamiltonian
of the preceding notes. Let $\mathcal H_m$ be the $m$-fermion sector,
$h=H_m$, and $H_-=H_{m-1}$. For $m>0$, define
\begin{equation}
 V:\mathcal H_m\longrightarrow\mathbb C^N\otimes\mathcal H_{m-1},
 \qquad V\psi=\frac1{\sqrt m}\sum_i|i\rangle c_i\psi,
 \qquad K=1_N\otimes H_- .
\end{equation}
Since $\sum_i c_i^\dagger c_i=m$, $V^\dagger V=1$.
The unital completely positive map $\Phi(Z)=V^\dagger(1\otimes Z)V$
is explicit; no emergent spatial coordinate has been assumed.
For $S(t_a,t_b)=\sum_i c_i^\dagger(t_a)c_i(t_b)$, set
\begin{equation}\label{eq:W}
 W(t)=e^{iKt}Ve^{-iht},\qquad
 B_{ab}=\frac{S_m(t_a,t_b)}m=W(t_a)^\dagger W(t_b).
\end{equation}
Thus $B_{aa}=1$, $B_{ab}^\dagger=B_{ba}$, and $\|B_{ab}\|\leq1$.
For arbitrary vectors $\psi_a$, positivity is exact:
\begin{equation}
 \sum_{a,b}\langle\psi_a,B_{ab}\psi_b\rangle
 =\left\|\sum_bW(t_b)\psi_b\right\|^2\geq0.
\end{equation}
This is a concrete representation of the entire two-time object,
including the same-sample correlation between neighbouring charge
Hamiltonians. A generic independent-probe kernel would lose it.

\newpage
\section{Composition and the full projected algebra}
Let $P_b=W(t_b)W(t_b)^\dagger$. The common-middle-time product is
\begin{equation}\label{eq:composition}
 B_{ab}B_{bd}=B_{ad}-W(t_a)^\dagger(1-P_b)W(t_d).
\end{equation}
In particular $1-B_{ab}B_{ba}\geq0$. The defect describes amplitude
outside the charge-removal image. For arbitrary times the product is
$B_{ab}B_{cd}=W_a^\dagger W_bW_c^\dagger W_d$; subtracting the
reversed product gives the exact projected commutator. This expression
keeps the operator insertions that were left unresolved by the previous
unprojected CAR-norm calculation.

There is also a useful energy-space expression. In the eigenbasis of
$h$ and $H_-$ write $c^i_{\alpha a}=\langle\alpha,m-1|c_i|a,m\rangle$.
For the centred bilocal $S_c(\delta)=S(-\delta/2,\delta/2)$,
\begin{equation}\label{eq:spectral}
 (S_c(\delta))_{ab}
 =\sum_{i,\alpha}(c^i_{\alpha a})^*c^i_{\alpha b}
 e^{i\delta[E_\alpha-(E_a+E_b)/2]}.
\end{equation}
Consequently its covariance is a positive Gram kernel in $\delta$,
obtained by summing $S_c(\delta)_{ab}^*S_c(\epsilon)_{ab}$ over $a,b$.
The intermediate amplitudes must be added coherently before taking
absolute squares. Their internal frequency is an adjacent-sector
energy combination, not yet a quark momentum fraction or meson mass.

\section{The exact memory equation}
Put $\Pc=VV^\dagger$, $\Qc=1-\Pc$ and
\begin{equation}
 h_{\rm eff}=V^\dagger KV=\left(1-\frac pm\right)h,
 \qquad R=\Qc KV,\qquad K_\perp=\Qc K\Qc.
\end{equation}
The effective Hamiltonian follows from
$\sum_i c_i^\dagger Hc_i=(m-p)h$. Resolve $e^{iKt}V$ into its
image and orthogonal components. The latter has zero initial value;
integrating it and substituting back yields
\begin{align}\label{eq:memory}
 F(t)&=V^\dagger e^{iKt}V,\qquad F(0)=1,\\
 \dot F(t)&=ih_{\rm eff}F(t)-\int_0^t\mathcal M(t-s)F(s)\,ds,
 &\mathcal M(s)&=R^\dagger e^{iK_\perp s}R.\nonumber
\end{align}
This is the standard projection-memory construction \cite{Mori},
here applied to explicit finite matrices. No infinite-dimensional
domain assumption is required. The associated Schur complement is
\begin{equation}
 V^\dagger(z-K)^{-1}V=
 [z-h_{\rm eff}-R^\dagger(z-K_\perp)^{-1}R]^{-1},
 \quad\operatorname{Im}z\ne0.
\end{equation}

\newpage
\section{The existing singlet is the initial memory fluctuation}
Let $A=\sum_i[H,c_i]^\dagger[H,c_i]$. The known CAR identity
$\sum_i c_i^\dagger H^2c_i=(m-2p)h^2+A_m$ gives
\begin{equation}\label{eq:G}
 G:=\mathcal M(0)=R^\dagger R
 =\frac{A_m}{m}-\frac{p^2}{m^2}h^2\geq0.
\end{equation}
In particular $A_m\geq p^2h^2/m$ and
$[h,G]=[h,A_m]/m$. For $m=p$, $H_-=0$, $G=0$ and the
compression is exactly autonomous. For general interacting sectors
it has nonzero memory. Even the free $p=1$ model can have $G\ne0$:
memory of a compression does not by itself diagnose interacting mesons.

At finite $N$, expansion of the full bilocal gives
\begin{equation}
 S_c(\delta)=m-ip\delta h-\tfrac12\delta^2 A_m+O(\delta^3),
 \qquad
 \frac{S_c(\delta)+S_c(-\delta)}2
 =m-\tfrac12\delta^2 A_m+O(\delta^4).
\end{equation}
Thus the first nonconserved even-relative-time coefficient is exactly
the previously studied $A$, after the full flavour sum. At two centre
times, the coefficient of $\delta^2\epsilon^2$ in the commutator of
the even bilocals is $[A(T),A(T')]/4$. This is a local expansion,
not an assertion of uniformity of the derivative expansion in $p$.

To connect to the existing dynamical channel, on each sector use
$C_m=\EE\trm h^2$, $a_m=\EE\trm A_m$, and the earlier ensemble
projection coefficient $c_m$ onto $h^2-C_m$. Define
\begin{equation}
 Y_m=A_m-a_m1-c_m(h^2-C_m1).
\end{equation}
Equation \eqref{eq:G} becomes the exact decomposition
\begin{equation}\label{eq:decomp}
 G=\left(\frac{a_m}m-\frac{p^2C_m}{m^2}\right)1
 +\left(\frac{c_m}m-\frac{p^2}{m^2}\right)(h^2-C_m1)
 +\frac{Y_m}m.
\end{equation}
Applying this sector by sector, $mG_{\rm residual}=Y$ is precisely
the earlier singlet after its specified conserved subtraction.
In the full annealed trace its DS covariance is already known:
\begin{equation}
 \EE\tr_0[Y(T)Y(0)]\longrightarrow W_dG_{q^2}(T),\quad
 W_d=4J_s^4\left[\lambda^{-1}+\frac{4q}{1+q}\right],\quad
 q=e^{-4\lambda}.
\end{equation}
This is an external centre-time correlation of the initial memory
operator. It is \emph{not} the orthogonal-time memory function
$\mathcal M(s)$ in \eqref{eq:memory}; computing that function is the
next substantive problem. The distinction prevents replacing the
unknown dynamics by a previously known two-point function.

\newpage
\section{A controlled finite-time double-scaling statement}
Use the earlier Gaussian model with standard-CAR coefficient variance
\begin{equation}
 w=4^p\frac{J_s^2N}{p^2}\left(\frac{p!}{N^p}\right)^2,
 \qquad C_m=w\binom mp\binom{N-m+p}{p}.
\end{equation}
The sector second moment is an existing embedded-ensemble formula
\cite{Tall}. Taking a trace of the CAR identity gives
$a_m=(2p-m)C_m+mC_{m-1}$. Since
\begin{equation}
 \frac{C_{m-1}}{C_m}=\frac{m-p}{m}\frac{N-m+p+1}{N-m+1},
\end{equation}
the initial memory has the positive exact mean
\begin{equation}\label{eq:mean}
 \EE\trm G=C_m\frac{p(m-p)(N+1)}{m^2(N-m+1)}.
\end{equation}
For $m<p$, both $h$ and $G$ vanish. At half filling and
$p^2/N\to\lambda>0$, $C_m\to J_s^2/\lambda$ and
\begin{equation}
 p\,\EE\trm G\longrightarrow4J_s^2.
\end{equation}
For $N$ even and $m=N/2$, the first limit follows explicitly from
\begin{equation}
 C_{N/2}=\frac{J_s^2}{\lambda_N}
 \prod_{j=0}^{p-1}\left(1-\frac{2j}N\right)
 \prod_{j=1}^{p}\left(1+\frac{2j}N\right),\qquad
 \lambda_N=p^2/N.
\end{equation}
The leading logarithms cancel; the remainder is
$O(p/N+p^3/N^2)$.

More than a curvature limit is available. The intertwining defect
$KV-Vh_{\rm eff}=R$ gives the exact Duhamel formula
\begin{equation}
 e^{iKt}V-Ve^{ih_{\rm eff}t}
 =i\int_0^t e^{iK(t-s)}R e^{ih_{\rm eff}s}\,ds.
\end{equation}
Using the normalized Hilbert--Schmidt norm on the source sector,
unitarity and $V^\dagger$ contractivity imply
\begin{align}
 \|F(t)-e^{ih_{\rm eff}t}\|_{2,m}^2&\leq t^2\trm G,\\
 \EE\|B_{0t}-1\|_{2,m}^2
 &\leq2t^2\left[\EE\trm G+\frac{p^2}{m^2}C_m\right].\label{eq:bound}
\end{align}
The right-hand side is $O(J_s^2t^2/p)$ at half filling. It controls
the supremum over any fixed compact time interval. It does not extend
to times growing like $\sqrt p$ or to a rescaled fluctuation field
without additional estimates.

\newpage
\section{Full-trace limit and the reconstruction variable}
The same conclusion holds for the original infinite-temperature full
trace. Set $B_m=S_m/m$ for $m>0$ and $B_0=1$. Charge weights are
$2^{-N}\binom Nm$. In the central window
$|m-N/2|\leq N^{2/3}$, the product expansions above are uniform:
\begin{equation}
 C_m\to J_s^2/\lambda,\qquad
 p\,\EE\trm G\to4J_s^2,\qquad (p/m)^2C_m=O(p^{-2}).
\end{equation}
For example, the extra logarithmic errors are bounded by constants
times $p(m-N/2)^2/N^2+p^2|m-N/2|/N^2$, both tending to zero.
Outside this window the binomial probability is at most
$2e^{-2N^{1/3}}$. Uniform polynomial bounds suffice there:
\begin{equation}
 C_m\leq\frac{J_s^2}{\lambda_N}(1+p/N)^{2p}
 \leq\frac{J_s^2}{\lambda_N}e^{2\lambda_N},\qquad
 \EE\trm G\leq(N+1)C_m.
\end{equation}
The first inequality bounds the falling and rising products by
$[m(N-m+p)]^p\leq[(N+p)^2/4]^p$.
Hence averaging \eqref{eq:bound} gives
$\EE\|B_{0t}-1\|_{2,0}^2=O(t^2/p)$.
For the conventional normalized singlet,
\begin{equation}
 \EE\left\|\frac{2S(0,t)}N-1\right\|_{2,0}^2
 \leq8\EE\|B_{0t}-1\|_{2,0}^2+\frac2N
 \longrightarrow0.
\end{equation}
This uses $0\leq2Q/N\leq2$ and
$\tr_0(2Q/N-1)^2=1/N$. Centre-time conjugation preserves the norm,
so the statement is uniform when $|t_b-t_a|$ is bounded.

The ordinary normalized bilocal therefore has a trivial leading
limit, while the conserved-subtracted even derivative gives the
nontrivial $Y$ field in \eqref{eq:decomp}. A meson map must retain
such fluctuations. This establishes a normalization requirement,
not the complete fluctuation limit at finite relative time.

The polarization question also has useful existing guidance.
Xu's tracial chord representation \cite{Xu} separates field operators
from auxiliary creation/annihilation operators. A nonzero bounded
annihilator of a separating trace vector cannot belong to that same
observable algebra. Rajgadia--Xu \cite{RX} construct state-adapted
observables in a different KM-state setting. These are concrete
routes to examine, but their state and Hamiltonian must be matched
before transferring them to this charge-conserving model.

\newpage
\section*{Validation and the next calculation}
Four deterministic Hamiltonians verify twenty Gaussian-integer matrix
identities, including the positive factorization
\begin{equation}
 Z^\dagger Z=m^2 A_m-mp^2h^2,\qquad
 Z=mK\mathsf C-(m-p)\mathsf C h,\qquad \mathsf C=\sqrt m V.
\end{equation}
The same systems test the full bilocal against the dilation, positive
block kernels, the composition defect, the resolvent Schur complement,
24/48-node memory integration, and twelve finite-time bounds.
Five exact monomial covariance sums verify \eqref{eq:mean}.
Eighteen 80-digit evaluations retain finite-$p$ deviations. For example,
at $p=256$, $J_s=1$, the values of $p\EE\tr_mG$ are
$3.99867524$, $3.99473404$, $3.97903814$ for
$\lambda=0.25,0.5,1$.

An independent script checks the complete four-anticommutator formula
at unequal times in two interacting sectors, and reconstructs the
centred bilocal from coherent charged eigenstate amplitudes.
Even-relative-time difference quotients converge quadratically to
$A_m$. The four contributions individually have squared norms much
larger than their summed commutator, explicitly verifying the need
to keep cancellations. These finite-system examples do not establish
a nonzero DS limit of that commutator.

The next task is to calculate the covariance and spectral structure of
$R^\dagger e^{iK_\perp s}R$ after the required conserved subtraction
and normalization. This retains the explicit neighbouring-charge
correlation and supplies an actual kernel to compare with bulk
internal dynamics. Its memory parameter $s$, the external centre time,
and the light-front quark fraction must remain distinct until a
reconstruction is derived. No bulk mass has been assigned by fiat.

The dilation, Schur complement and projection-memory methods are
standard mathematics. The present result is their application to the
specified complex SYK bilocal; no priority claim is made. Numerical
checks use ordinary floating point where indicated, not certified
interval arithmetic. Analytic bounds support the finite-time limit.

\begin{thebibliography}{9}
\bibitem{Mori} C.~Widder, J.~Zimmer and T.~Schilling,
\emph{On the generalized Langevin equation and the Mori projection
operator technique}, \href{https://arxiv.org/abs/2503.20457}{arXiv:2503.20457}.
\bibitem{Tall} J.~Tall and S.~Tomsovic,
\emph{Double-scaled bosonic and fermionic embedded ensembles,
complex SYK, and the dual Hilbert space},
\href{https://arxiv.org/abs/2604.14522v2}{arXiv:2604.14522v2}.
\bibitem{Xu} J.~Xu, \emph{Von Neumann Algebras in Double-Scaled SYK},
\href{https://arxiv.org/abs/2403.09021v7}{arXiv:2403.09021v7}, Sec.~3.3.
\bibitem{RX} H.~Rajgadia and J.~Xu,
\emph{Emergent States and Algebras from the Double-Scaling limit of
Pure States in SYK}, \href{https://arxiv.org/abs/2604.14387v3}{arXiv:2604.14387v3},
Sec.~3.2.
\end{thebibliography}
\end{document}
